\documentclass[sigconf,nonacm]{acmart}

\settopmatter{printacmref=false}
\renewcommand\footnotetextcopyrightpermission[1]{}
\pagestyle{plain}

\usepackage{amsmath}
\usepackage{booktabs}
\usepackage{array}
\usepackage{enumitem}
\usepackage{hyperref}

\title{A Monte Carlo Simulation of Player Retention Under Alternative Pity Systems in an FGO-Style Gacha Model}

\author{Your Name}
\affiliation{%
  \institution{Your University}
  \city{Makati}
  \country{Philippines}
}
\email{your.email@example.com}

\begin{document}

\begin{abstract}
Gacha systems are widely used in free-to-play games as randomized reward mechanisms where players spend in-game or real-money currency for a chance to obtain rare virtual items or characters. Although pity systems are often introduced to reduce extreme unlucky outcomes, not all pity designs provide the same level of fairness, perceived progress, or player retention. This study proposes a Monte Carlo simulation of an FGO-style gacha model to compare the effects of hard pity, improved drop rates, soft pity, carryover pity, and token-based exchange systems on target acquisition rate, average pulls required, extreme bad-luck frequency, frustration score, and simulated player retention. The study uses synthetic player profiles such as free-to-play, low spender, medium spender, and high spender segments to evaluate how different gacha designs affect different resource levels. By modeling repeated summon attempts across multiple banners, the study aims to identify which gacha improvement produces the best balance between fairness and long-term retention. The expected contribution of this research is a practical and reproducible simulation framework for evaluating randomized monetization systems without requiring private player data.
\end{abstract}

\begin{CCSXML}
<ccs2012>
 <concept>
  <concept_id>10010405.10010489.10010495</concept_id>
  <concept_desc>Applied computing~Computer games</concept_desc>
  <concept_significance>500</concept_significance>
 </concept>
 <concept>
  <concept_id>10002950.10003648.10003671</concept_id>
  <concept_desc>Mathematics of computing~Simulation evaluation</concept_desc>
  <concept_significance>500</concept_significance>
 </concept>
 <concept>
  <concept_id>10003120.10003121</concept_id>
  <concept_desc>Human-centered computing~Human computer interaction (HCI)</concept_desc>
  <concept_significance>300</concept_significance>
 </concept>
</ccs2012>
\end{CCSXML}

\ccsdesc[500]{Applied computing~Computer games}
\ccsdesc[500]{Mathematics of computing~Simulation evaluation}
\ccsdesc[300]{Human-centered computing~Human computer interaction (HCI)}

\keywords{Monte Carlo simulation, gacha system, pity system, loot boxes, player retention, free-to-play games, randomized rewards, game monetization, simulation modeling}

\maketitle

\section{Introduction}

Free-to-play games have become one of the dominant business models in digital gaming, especially in mobile and live-service environments. In this model, games are offered without an initial purchase price, while revenue is generated through optional in-game purchases, cosmetics, premium currencies, battle passes, loot boxes, or gacha-based systems. Because revenue depends on continued participation rather than one-time purchase, player retention becomes a central concern for game developers and researchers. Musta\v{c} et al. emphasized that studying and modeling churn is highly important in free-to-play games because player engagement is closely connected to long-term revenue and sustainability \cite{mustac2022churn}. Similarly, Burelli noted that free-to-play games create unique modeling challenges because player behavior and spending patterns are often irregular and difficult to predict \cite{burelli2022clv}.

Among free-to-play monetization mechanics, gacha systems are particularly notable because they rely on randomized rewards. In a gacha system, players spend a resource, such as premium currency, tickets, or real money, to receive a random item or character. The desired reward is usually assigned a low probability, which means players may need many attempts before obtaining the target. These mechanics are closely related to loot boxes, which have been studied in relation to probability disclosure, spending behavior, player harm, and gambling-like design features \cite{xiao2022pandora,xiao2022probability,gibson2022microtransactions,xiao2023odds}. However, gacha systems are not uniform. Ballou et al. argued that loot-box-like mechanics vary widely across games and identified multiple design features that may influence player behavior or spending \cite{ballou2022hidden}. This suggests that gacha systems should not be treated as one generic mechanic; instead, individual design choices such as drop rates, pity thresholds, carryover rules, and exchange systems should be evaluated separately.

One common design intervention in gacha systems is the pity system. A pity system guarantees or increases the chance of obtaining a rare reward after a player has failed a certain number of times. Gan's economic model of gacha pricing describes this type of mechanism as a form of worst-case insurance after repeated failures \cite{gan2022gacha}. In practical game design, pity systems can reduce the possibility of extreme unlucky outcomes, but their effectiveness depends on the specific rules used. For example, a hard pity system may guarantee the target only after a very large number of pulls, while a soft pity system gradually increases the probability of success after repeated failures. A carryover pity system allows unsuccessful progress to transfer to future banners, while a token or spark system gives players a separate currency that can eventually be exchanged for the target reward.

This study uses an FGO-style gacha model as the baseline because \textit{Fate/Grand Order} is a recognizable example of a character-based gacha game with a hard pity structure. Public game references describe its rate-up SSR model as having a low target acquisition probability and a pity guarantee at a high pull threshold \cite{fgowiki2024pity}. This makes it useful as a simulation case because the system provides a clear example of a low-probability target reward with a late guarantee. However, the study does not claim to evaluate the actual commercial performance of \textit{Fate/Grand Order} or any specific player community. Instead, the term ``FGO-style'' refers only to a simplified gacha structure with low target probability, limited banners, and a hard pity rule.

Recent literature shows that probability disclosures and pity timers are relevant to consumer protection and player perception. Xiao, Fraser, and Newall studied player opinions on probability disclosures and pity timers in China and found that many players were aware of pity timers and considered them appropriate \cite{xiao2022pandora}. However, Xiao and Newall also argued that probability disclosures alone may not be enough because loot-box reward structures can be complex and difficult for players to evaluate \cite{xiao2022probability}. Xiao, Henderson, and Newall further found poor compliance with probability disclosure self-regulation in the United Kingdom market \cite{xiao2023odds}. These studies support the need to examine not only whether probabilities are disclosed, but also whether the structure of the gacha system itself creates more understandable, fair, and retention-supportive outcomes.

Despite the growing body of research on loot boxes, microtransactions, and problematic spending, fewer studies directly simulate how alternative gacha pity systems affect player outcomes across repeated banners. Much of the existing literature focuses on associations between loot-box purchasing and gambling-related risk, financial harm, or mental wellbeing \cite{gibson2022microtransactions,sirola2023indebtedness,gonzalez2024lootboxes}. These studies are valuable for establishing the importance of randomized reward systems, but they do not directly answer a design question that is important for modeling and simulation: if a gacha system must exist, which pity design produces better fairness and retention outcomes under controlled assumptions?

This study addresses that gap by using Monte Carlo simulation to compare alternative gacha designs. Instead of collecting private player data, the study generates synthetic players with different pull budgets and behavioral assumptions. Each simulated player attempts to obtain a target reward across multiple banners. The simulation records whether the player obtains the target, how many pulls are required, whether pity is reached, whether progress is wasted, and whether the player remains active after repeated success or failure. Through this approach, the study can compare gacha systems using measurable outputs such as target acquisition rate, average pulls required, frustration score, and retention probability.

\section{Background of the Study}

Gacha systems are based on random probability. If the probability of obtaining a target reward in one pull is represented as $p$, then the probability of failing to obtain the target after $n$ independent pulls is:

\begin{equation}
P(\text{failure after } n \text{ pulls}) = (1-p)^n.
\label{eq:failure}
\end{equation}

Therefore, the probability of obtaining at least one target reward after $n$ pulls is:

\begin{equation}
P(\text{success within } n \text{ pulls}) = 1 - (1-p)^n.
\label{eq:success}
\end{equation}

These equations show why low-probability gacha systems can create large variation in player experience. Two players may spend the same amount of currency but receive very different outcomes. One player may obtain the target early, while another may fail after hundreds of pulls. This variation creates the need for pity systems, which limit the worst-case outcome by guaranteeing the target or increasing its probability after a certain number of failed attempts.

However, a pity system does not automatically guarantee strong retention. A pity threshold may be so high that only high-spending or long-saving players can realistically reach it. For free-to-play or low-spending players, a non-carryover pity system may still feel punishing because partial progress disappears when a banner ends. In this case, the player may experience a sense of wasted effort even if the system technically has a guarantee. This is important because retention is not only affected by mathematical probability, but also by perceived fairness, frustration, and the feeling of progress.

The study therefore models retention as an outcome affected by gacha experience. A simplified retention model may treat successful pulls as positive experiences and repeated failures, wasted pity progress, and currency depletion as negative experiences. For example, a player's frustration score may increase when the player spends many pulls without receiving the target, fails several banners in a row, or loses pity progress after a banner ends. Retention probability can then be estimated as a decreasing function of frustration and an increasing function of successful acquisition or preserved progress.

A possible frustration score can be represented as:

\begin{equation}
F_i = \alpha L_i + \beta B_i + \gamma W_i - \delta S_i,
\label{eq:frustration}
\end{equation}

where $F_i$ is the frustration score of player $i$, $L_i$ is the number of failed pulls, $B_i$ is the number of failed banners, $W_i$ is the amount of wasted pity progress, $S_i$ is the number of successful target acquisitions, and $\alpha$, $\beta$, $\gamma$, and $\delta$ are adjustable weights.

The simulated retention probability may then be modeled as:

\begin{equation}
R_i = \max(0, \min(1, 1 - \lambda F_i)),
\label{eq:retention}
\end{equation}

where $R_i$ is the retention probability of player $i$ and $\lambda$ is a scaling parameter that controls how strongly frustration affects retention. This equation is intentionally simple because the study focuses on comparing gacha system designs under consistent assumptions rather than predicting exact real-world retention rates.

\section{Problem Statement}

Although pity systems are commonly used to reduce extreme bad luck in gacha games, their actual effect on player retention may vary depending on the design. A hard pity system with a high threshold may protect only the unluckiest or highest-resource players. A simple drop-rate increase may improve average success but may not solve the issue of wasted progress. A carryover pity system or token-based system may improve perceived fairness because unsuccessful pulls still contribute to future acquisition. However, these alternatives must be evaluated using comparable metrics.

The problem addressed by this study is the lack of a simple, reproducible simulation-based comparison of alternative pity systems in an FGO-style gacha model. Specifically, the study investigates whether improving the base rate, adding soft pity, allowing pity carryover, or introducing a token/spark system produces better simulated retention and fairness outcomes than a baseline hard-pity-only system.

\section{Objectives of the Study}

The general objective of this study is to simulate and compare the effect of alternative gacha pity systems on target acquisition, player frustration, and retention in an FGO-style gacha model.

The specific objectives are as follows:

\begin{enumerate}[label=\textbf{O\arabic*.}]
    \item To construct a baseline FGO-style gacha simulation using a low target SSR rate and a hard pity threshold.
    \item To compare alternative gacha improvement systems, including improved base rate, soft pity, carryover pity, and token-based exchange.
    \item To evaluate each system using target acquisition rate, average pulls required, pity activation rate, extreme bad-luck rate, frustration score, and simulated retention rate.
    \item To compare outcomes across different synthetic player segments, such as free-to-play, low spender, medium spender, and high spender profiles.
    \item To identify which gacha improvement provides the best balance between player fairness and simulated long-term retention.
\end{enumerate}

\section{Research Questions}

This study aims to answer the following research questions:

\begin{enumerate}[label=\textbf{RQ\arabic*.}]
    \item How does a baseline hard-pity-only gacha system perform in terms of target acquisition rate, average pulls required, and extreme bad-luck frequency?
    \item How do improved base rates, soft pity, carryover pity, and token-based exchange systems affect simulated player retention compared with the baseline system?
    \item Which pity system provides the greatest retention improvement for free-to-play and low-spending players?
    \item Does preserving failed progress through carryover pity or token exchange reduce simulated frustration more effectively than simply increasing the target drop rate?
    \item Which gacha design produces the best balance between fairness, accessibility, and retention across multiple player segments?
\end{enumerate}

\section{Scope and Limitations}

This study is limited to a simulation-based analysis of gacha mechanics. It does not collect real spending data, private player logs, or survey responses from actual players. Instead, the study uses synthetic player profiles and assumed behavioral parameters. The simulation focuses on obtaining one target SSR or rare character per banner and does not model additional copies, character power levels, competitive rankings, story attachment, event rewards, or social influence.

The study also does not claim to represent the full commercial design of \textit{Fate/Grand Order} or any specific game. The FGO-style baseline is used only as a simplified model of a low-rate, hard-pity gacha system. The actual experience of players may differ due to emotional attachment to characters, prior savings, limited-time banners, personal spending limits, community influence, and other psychological factors.

Another limitation is that retention is modeled indirectly. Since actual retention data is not used, retention probability is estimated using a frustration-based formula. This means that the simulation can compare relative outcomes between gacha systems, but it cannot claim exact real-world retention rates. The results should therefore be interpreted as exploratory and comparative rather than predictive of actual player behavior.

\section{Significance of the Study}

This study is significant for several groups.

For game designers, the study provides a simulation-based method for testing whether a gacha system is likely to feel fair across different player segments. Instead of relying only on average revenue or target probability, designers can evaluate frustration, wasted progress, and retention-related outcomes.

For players, the study can help explain why some gacha systems feel more punishing than others even when they have similar published drop rates. It can also show how mechanics such as carryover pity and token exchange reduce wasted progress.

For researchers, the study contributes a reproducible simulation framework for analyzing randomized reward systems. It connects gacha probability modeling with player retention, which is an important concern in free-to-play game analytics.

For regulators and consumer protection discussions, the study may provide insight into why probability disclosure alone may not be enough. Even when drop rates are visible, the structure of pity systems, progress loss, and limited-time banners may still affect player behavior and perceived fairness.

\section{Proposed Simulation Framework}

The proposed simulation framework has three main components: input parameters, simulation process, and output metrics.

The input parameters include the base target rate, pity threshold, number of banners, monthly pull budget, player segment, carryover rule, soft pity rule, and token exchange threshold. Player segments are defined according to available pulls per banner or per month. For example, free-to-play players may have a lower pull budget, while high spenders may have enough pulls to reach pity more frequently.

The simulation process follows a repeated Monte Carlo structure. For each player and banner, the model randomly generates summon outcomes based on the assigned probability rules. If the player obtains the target, the simulation records the number of pulls used and updates the player's satisfaction or retention state. If the player fails, the model records remaining currency, failed progress, and frustration. In systems with carryover pity or tokens, unused progress is transferred to future banners. In systems without carryover, progress is reset.

\begin{table}[h]
\centering
\caption{Proposed gacha system scenarios for simulation.}
\label{tab:scenarios}
\begin{tabular}{p{0.18\linewidth}p{0.21\linewidth}p{0.43\linewidth}}
\toprule
\textbf{Scenario} & \textbf{Target Rate} & \textbf{Pity Rule} \\
\midrule
A: Baseline & 0.8\% & Hard pity at 330 pulls \\
B: Improved Rate & 1.0\% to 1.2\% & Hard pity at 330 pulls \\
C: Soft Pity & 0.8\% base & Rate increases after a defined failed-pull threshold \\
D: Carryover Pity & 0.8\% base & Failed pity progress transfers across banners \\
E: Token/Spark & 0.8\% base & Every pull gives tokens; target can be exchanged after threshold \\
\bottomrule
\end{tabular}
\end{table}

The output metrics include target acquisition rate, average pulls required to obtain the target, median pulls required, percentage of players reaching pity, percentage of players failing after spending all available pulls, extreme bad-luck rate, average frustration score, simulated retention rate, and segment-level retention difference between free-to-play, low spender, medium spender, and high spender players.

\section{Expected Contribution}

The expected contribution of this study is a practical Monte Carlo simulation model for comparing gacha pity systems. The study does not attempt to prove the real-world psychological effects of gacha mechanics. Instead, it provides a controlled way to compare design alternatives using consistent assumptions. The expected finding is that systems preserving player progress, such as carryover pity or token exchange, may produce better retention-related outcomes than systems that only increase the base drop rate. This is because preserved progress may reduce the feeling of wasted effort, especially for free-to-play and low-spending players.

\begin{thebibliography}{99}

\bibitem{gan2022gacha}
T. Gan, ``Gacha Game: When Prospect Theory Meets Optimal Pricing,'' \textit{arXiv preprint arXiv:2208.03602}, 2022. [Online]. Available: https://arxiv.org/abs/2208.03602

\bibitem{xiao2022pandora}
L. Y. Xiao, T. C. Fraser, and P. W. S. Newall, ``Opening Pandora's Loot Box: Novel Links with Gambling, and Player Opinions on Probability Disclosures and Pity-Timers in China,'' \textit{Journal of Gambling Studies}, vol. 39, pp. 645--668, 2022. doi: 10.1007/s10899-022-10148-0.

\bibitem{xiao2022probability}
L. Y. Xiao and P. W. S. Newall, ``Probability Disclosures Are Not Enough: Reducing Loot Box Reward Complexity as a Part of Ethical Video Game Design,'' \textit{Journal of Gambling Issues}, 2022. doi: 10.4309/jgi.2022.49.7.

\bibitem{ballou2022hidden}
N. Ballou, C. Gbadamosi, and D. Zendle, ``The Hidden Intricacy of Loot Box Design: A Granular Description of Random Monetized Reward Features,'' in \textit{Proceedings of DiGRA 2022 Conference: Bringing Worlds Together}, 2022. doi: 10.26503/dl.v2022i1.1327.

\bibitem{mustac2022churn}
K. Musta\v{c}, K. Ba\v{c}i\'{c}, L. Skorin-Kapov, and M. Su\v{z}njevi\'{c}, ``Predicting Player Churn of a Free-to-Play Mobile Video Game Using Supervised Machine Learning,'' \textit{Applied Sciences}, vol. 12, no. 6, article 2795, 2022. doi: 10.3390/app12062795.

\bibitem{burelli2022clv}
P. Burelli, ``Predicting Customer Lifetime Value in Free-to-Play Games,'' \textit{arXiv preprint arXiv:2209.12619}, 2022. [Online]. Available: https://arxiv.org/abs/2209.12619

\bibitem{gibson2022microtransactions}
E. Gibson, M. D. Griffiths, F. Calado, and A. Harris, ``The Relationship Between Videogame Micro-Transactions and Problem Gaming and Gambling: A Systematic Review,'' \textit{Computers in Human Behavior}, vol. 131, article 107219, 2022. doi: 10.1016/j.chb.2022.107219.

\bibitem{xiao2023odds}
L. Y. Xiao, L. L. Henderson, and P. W. S. Newall, ``What Are the Odds? Poor Compliance with UK Loot Box Probability Disclosure Industry Self-Regulation,'' \textit{PLOS ONE}, vol. 18, no. 9, e0286681, 2023. doi: 10.1371/journal.pone.0286681.

\bibitem{sirola2023indebtedness}
A. Sirola, R. Han, and A. Strzelecki, ``Loot Box Purchasing and Indebtedness: The Role of Psychosocial Factors and Problem Gambling,'' \textit{Addictive Behaviors Reports}, 2023.

\bibitem{gonzalez2024lootboxes}
J. Gonz\'{a}lez-Cabrera, V. Caba-Machado, A. D\'{i}az-L\'{o}pez, and colleagues, ``The Mediating Role of Problematic Use of Loot Boxes Between Internet Gaming Disorder and Online Gambling Disorder: Cross-Sectional Analytical Study,'' \textit{JMIR Serious Games}, vol. 12, e57304, 2024. doi: 10.2196/57304.

\bibitem{fgowiki2024pity}
Fate/Grand Order Wiki, ``Rate-Up SSR Pity System,'' 2024. Accessed: Jun. 29, 2026. [Online]. Available: https://fategrandorder.fandom.com/wiki/Rate-Up_SSR_Pity_System

\end{thebibliography}

\end{document}
